\documentclass{ximera}
%% Macros shared by every activity. Each activity loads this file with
%% %% Macros shared by every activity. Each activity loads this file with
%% %% Macros shared by every activity. Each activity loads this file with
%% \input{../preamble.tex}, and the xourse file with \input{preamble.tex}.

%% MathML intent for screen readers: \xmunits, \xmabs, \xmpoint and the
%% other commands of TailorSwiftBot.sty. The file is latex/xmintent.sty.
\usepackage{xmintent}

%% The interactive pieces of Ximera: \answer, multipleChoice and \choice,
%% selectAll, hint and the rest. Keep this line: without it those commands
%% are undefined and come out as plain text. The file is
%% latex/ximera-components.sty.
\usepackage{ximera-components}

%% Tables with the tabularray package. The settings rowhead, xmcolhead,
%% caption and summary make their header cells known to screen readers, on
%% the web page and in the tagged PDF.
\usepackage{tabularray}

\newcommand{\RR}{\mathbb{R}}
\newcommand{\dfn}{\textit}
, and the xourse file with \input{preamble.tex}.

%% MathML intent for screen readers: \xmunits, \xmabs, \xmpoint and the
%% other commands of TailorSwiftBot.sty. The file is latex/xmintent.sty.
\usepackage{xmintent}

%% The interactive pieces of Ximera: \answer, multipleChoice and \choice,
%% selectAll, hint and the rest. Keep this line: without it those commands
%% are undefined and come out as plain text. The file is
%% latex/ximera-components.sty.
\usepackage{ximera-components}

%% Tables with the tabularray package. The settings rowhead, xmcolhead,
%% caption and summary make their header cells known to screen readers, on
%% the web page and in the tagged PDF.
\usepackage{tabularray}

\newcommand{\RR}{\mathbb{R}}
\newcommand{\dfn}{\textit}
, and the xourse file with \input{preamble.tex}.

%% MathML intent for screen readers: \xmunits, \xmabs, \xmpoint and the
%% other commands of TailorSwiftBot.sty. The file is latex/xmintent.sty.
\usepackage{xmintent}

%% The interactive pieces of Ximera: \answer, multipleChoice and \choice,
%% selectAll, hint and the rest. Keep this line: without it those commands
%% are undefined and come out as plain text. The file is
%% latex/ximera-components.sty.
\usepackage{ximera-components}

%% Tables with the tabularray package. The settings rowhead, xmcolhead,
%% caption and summary make their header cells known to screen readers, on
%% the web page and in the tagged PDF.
\usepackage{tabularray}

\newcommand{\RR}{\mathbb{R}}
\newcommand{\dfn}{\textit}


\title{Accessibility Showcase}

\begin{document}
\begin{abstract}
One example of every accessibility feature of this sandbox, on one page.
\end{abstract}
\maketitle

This page exists to be tested. Read it with a screen reader as a web page,
then open its tagged PDF (the link at the top right) and read that, or
look at the Tags panel in Adobe Acrobat. Each part says what to listen or
look for.

\section{Headings}

The title of the page is a level 1 heading and each section a level 2
heading. A subsection is one level further down, and so is the head of a
theorem inside it.

\subsection{A subsection with a theorem}

\begin{theorem}[Pythagoras]
In a right triangle with legs $a$ and $b$ and hypotenuse $c$, we have
$a^2+b^2=c^2$.
\end{theorem}

\section{Formulas}

Each formula carries MathML, and these commands add what the notation
means, so that it is read correctly:

\begin{itemize}
\item an absolute value, $\xmabs{x-3}$, and a determinant, $\xmdet{A}$;
\item a length with a unit, $5\,\xmunits{cm}$, and an angle, $90\xmdeg$;
\item a point, $\xmpoint{2,5}$, and an interval, $\xminterval{[0,1)}$;
\item a composition, $\xmcompose{f}{g}(x)$, and an inverse, $\xminv{f}(x)$;
\item a dot product of vectors, $\xmdot{\xmvec{u}}{\xmvec{v}}$.
\end{itemize}

A displayed formula:
\[
\xmabs{a+b}\le\xmabs{a}+\xmabs{b}.
\]

\section{Blanks and answers}

A line to write on is announced as a blank. Name: \xmblank[2in]

\begin{problem}
An answer box. On the web page it is a box to type in; in the PDF it is a
blank. $2+2=\answer{4}$
\end{problem}

\section{Lists and choices}

A numbered list is tagged as a list, with its items:

\begin{enumerate}
\item Read the question.
\item Choose an answer.
\item Check it.
\end{enumerate}

\begin{problem}
On the web page these choices are radio buttons; in the PDF they are a
list. Which number is prime?

\begin{multipleChoice}
\choice{$21$}
\choice[correct]{$23$}
\choice{$27$}
\end{multipleChoice}
\end{problem}

\section{A table}

The first row holds column headers and the first column holds row headers.
Moving through the cells, a screen reader announces the headers of each one.

\begin{center}
\begin{tblr}{colspec={lcc}, hlines, vlines, rowhead=1, xmcolhead=1,
    caption={Quiz scores}, summary={One row for each student, one column for each quiz.}}
Student & Quiz 1 & Quiz 2 \\
Ada     & 9      & 10     \\
Emmy    & 8      & 7      \\
\end{tblr}
\end{center}

\section{Pictures}

A picture from a file, with a caption and a description. The description
is what a screen reader reads in place of the picture.

\begin{figure}
\begin{image}
  \includegraphics[width=3in]{sample.png}
\end{image}
\caption{A sample picture.}
\xmalt{A rectangle filled with a smooth colour gradient, blue on the left and orange on the right.}
\end{figure}

A drawn picture, with a description:

\begin{center}
\begin{tikzpicture}[scale=.75]
  \draw[<->] (-1,0)--(4,0) node[right]{$x$};
  \draw[<->] (0,-1)--(0,4) node[above]{$y$};
  \draw[thick] (-0.5,-0.5)--(3.5,3.5);
  \fill (2,2) circle (2pt) node[right]{$(2,2)$};
\end{tikzpicture}
\end{center}
\xmalt{Coordinate axes with the line y equals x drawn through the origin, and the point 2 comma 2 marked on it.}

\section{Words in another language}

The French phrase \xmlang{fr}{raison d'\^etre} is marked as French, so a
screen reader can switch voice for it and back.

The same works for other alphabets. The Ukrainian mathematician Mykhailo
Kravchuk wrote his name \xmlang{uk}{Михайло Кравчук}, and the Greek word
\xmlang{el}{τέλειος} means complete.

\section{Links and footnotes}

A link has a name that says where it goes:
\href{https://www.w3.org/WAI/fundamentals/accessibility-intro/}{Introduction to Web Accessibility, from the W3C}.

A footnote\footnote{On the web page the marker is a button that shows and
hides this note.} is tagged as a note.

\section{Text that opens and closes}

On the web page these are operated with a button that says whether the
text is open. In the PDF their text is simply printed.

\begin{hint}
A hint, closed until it is asked for.
\end{hint}

\begin{expandable}{xmexpandable}{A worked example}
To solve $2x+1=7$, subtract $1$ and divide by $2$.
\end{expandable}

\end{document}
